% !TeX root = main.tex
% Documentation: https://www.ctan.org/pkg/glossaries

% Create your glossary entries here

\newglossaryentry{latex}
{
    type=main,
    name=latex,
    description={Is a markup language specially suited
            for scientific documents, but this time with a description that is way too long to fit in just one line. I mean... what were i thinking making it so long? Ieserve this to fail but hopefully it wont.}
}

\newglossaryentry{maths}
{
    type=main,
    name=mathematics,
    description={Mathematics is what mathematicians do}
}

\newglossaryentry{Forschungsdaten}
{
    type=main,
    name=Forschungsdaten,
    description={
            \enquote{
                Zu Forschungsdaten zählen u.a. Messdaten, Laborwerte, audiovisuelle Informationen, Texte, Surveydaten, Objekte aus Sammlungen oder Proben, die in der wissenschaftlichen Arbeit entstehen, entwickelt oder ausgewertet werden. Methodische Testverfahren, wie Fragebögen, Software und Simulationen können ebenfalls zentrale Ergebnisse wissenschaftlicher Forschung darstellen und sollten daher ebenfalls unter den Begriff Forschungsdaten gefasst werden.
            } \cite{SENA15}
        }
}

\newglossaryentry{Forschungsdatenmanagement}
{
    name=Forschungsdatenmanagement,
    description={
            Beschreibt das Management von \Gls{Forschungsdaten}.
        }
}
\newglossaryentry{opc_ua}{
    type=main,
    name={OPC UA},
    description={Open Platform Communications Unified Architecture
        },
}
